\documentclass[11pt]{article}

% ----------------------------------------------------------------------------
% Urban Mobility Equity Dashboard --- MS Project Report, Initial Draft (Ch. 1--3)
% Styled to match the EMGT 7945 Research Project Proposal.
% Compile with: tectonic report/Chapters_1-3_Draft.tex
% ----------------------------------------------------------------------------

\usepackage[letterpaper,margin=1in]{geometry}
\usepackage{times}
\usepackage[T1]{fontenc}
\usepackage[utf8]{inputenc}
\usepackage{amsmath}
\usepackage{graphicx}
\usepackage{booktabs}
\usepackage{array}
\usepackage{tabularx}
\usepackage{enumitem}
\usepackage{titlesec}
\usepackage{caption}
\usepackage[hidelinks]{hyperref}
\usepackage{fancyhdr}
\usepackage{xcolor}

\definecolor{neunavy}{RGB}{21,51,79}

% Section styling to echo the proposal's numbered, bold uppercase headings.
\titleformat{\section}{\large\bfseries\color{neunavy}}{\thesection.}{0.5em}{\MakeUppercase}
\titleformat{\subsection}{\normalsize\bfseries}{\thesubsection}{0.5em}{}
\titlespacing*{\section}{0pt}{14pt}{6pt}
\titlespacing*{\subsection}{0pt}{8pt}{3pt}

\setlength{\parskip}{5pt}
\setlength{\parindent}{0pt}
\newcolumntype{Y}{>{\raggedright\arraybackslash}X}

\pagestyle{fancy}
\fancyhf{}
\rhead{\small Urban Mobility Equity Dashboard}
\lhead{\small EMGT 7945 --- Initial Draft (Ch. 1--3)}
\rfoot{\small\thepage}

\captionsetup[table]{labelfont=bf,font=small,skip=4pt}

\begin{document}

% ----------------------------- Title block -----------------------------
\thispagestyle{empty}
\begin{center}
{\large\bfseries MS PROJECT REPORT --- INITIAL DRAFT (CHAPTERS 1--3)}\\[2pt]
{\normalsize EMGT 7945 MS Project}\\[18pt]
{\LARGE\bfseries\color{neunavy} Urban Mobility Equity Dashboard:}\\[4pt]
{\Large\bfseries\color{neunavy} Mapping Who Gets Left Behind in U.S. Cities}\\[22pt]
\end{center}

\begin{center}
\renewcommand{\arraystretch}{1.3}
\begin{tabular}{r l}
\textbf{Submitted by} & Aashritha Krishna Kanagala \\
\textbf{Program} & Master of Science in Engineering Management (MSEM) \\
\textbf{Course} & EMGT 7945 MS Project \\
\textbf{Department} & Mechanical and Industrial Engineering \\
\textbf{University} & Northeastern University, Boston, Massachusetts \\
\textbf{Supervisor} & Professor Sivarit (Tony) Sultornsanee \\
\textbf{Date} & June 2026 \\
\end{tabular}
\end{center}

\vspace{14pt}
\noindent\fbox{\parbox{\dimexpr\textwidth-2\fboxsep-2\fboxrule}{\small
\textbf{Draft status.} This initial draft delivers Chapters 1--3 (Introduction,
Literature Review, Methodology) and reports the implemented pipeline, Mobility
Equity Score (MES), and interactive dashboard, all running end-to-end for the
three Phase-1 cities (Chicago, Houston, Seattle) on \emph{measured open data
across every dimension}: transit (GTFS --- CTA, METRO Houston, King County Metro),
walkability (EPA Smart Location Database), bicycle (OpenStreetMap), EV charging
(OpenStreetMap charging points), and demographics (Census ACS 2019--2023). One
data caveat carried to Chapter~5: the EV layer uses community-maintained OSM
charging points as a lower-bound proxy for the authoritative DOE/NREL AFDC
inventory. Chapter~4 (Results) is forthcoming; the analytical machinery is
validated on the live data.}}

\vspace{6pt}

% ============================== CHAPTER 1 ==============================
\section{Introduction and Motivation}

\subsection{1.1\quad Background}
Transportation is one of the most fundamental determinants of economic
opportunity in urban America. A person's ability to reach employment, healthcare,
education, and essential services depends on a layered ecosystem of
infrastructure: transit frequency and coverage, pedestrian networks, bicycle
facilities, and access to electric-vehicle (EV) charging. Each dimension has
documented equity gaps across U.S. cities. Transit deserts---areas where transit
supply falls far short of the demand of car-free, low-income residents---
disproportionately affect minority communities~\cite{aman2020}. Walkability is
highly racialized, with communities of color consistently inhabiting less
walkable neighborhoods~\cite{urban2022}. Disadvantaged communities have roughly
64\% fewer public EV chargers per capita than non-disadvantaged
communities~\cite{lou2025}, and bicycle infrastructure remains underfunded in
lower-income areas~\cite{ahmed2024}.

Despite this evidence, the available tools remain siloed. The EPA National
Walkability Index, the DOE Alternative Fuels Station Locator, the FTA National
Transit Database, and OpenStreetMap bicycle data each measure one dimension in
isolation. No open-access tool integrates all four into a single
neighborhood-level equity framework---the gap this project fills.

\subsection{1.2\quad Problem Statement}
Low-income and minority neighborhoods in major U.S. cities likely experience
compounded, multi-dimensional mobility disadvantage, scoring poorly not in one
dimension but across several simultaneously. The magnitude of this compounding
effect, and the specific neighborhoods where it is most severe, have not been
quantified in an integrated, publicly accessible format. This directly impedes
evidence-based infrastructure-investment decisions at the city and federal levels.

\subsection{1.3\quad Research Questions}
The project is guided by one central question and four sub-questions:
\begin{enumerate}[leftmargin=1.4em,itemsep=1pt]
\item \textbf{Central.} To what extent do low-income and minority neighborhoods
in major U.S. cities experience compounded disadvantage across transit access,
walkability, bicycle infrastructure, and EV charging, and how can this
multi-dimensional inequity be measured, visualized, and acted upon?
\item \textit{Measurement.} Can heterogeneous public datasets be integrated into
a single, transparent, tract-level composite index that is reproducible and
comparable across cities?
\item \textit{Distribution.} Where are the mobility deserts (bottom-quartile MES)
and compounded deserts (bottom quartile on $\geq 2$ dimensions) within each city?
\item \textit{Equity.} Is the MES associated with median income, racial
composition, vehicle access, and tenure, and in what direction and magnitude?
\item \textit{Robustness.} Are conclusions sensitive to the weighting of
sub-indices and to the percentile threshold defining a mobility desert?
\end{enumerate}

\subsection{1.4\quad Specific Objectives}
\begin{enumerate}[leftmargin=1.4em,itemsep=1pt]
\item Construct a Mobility Equity Score (MES) by aggregating four mobility
sub-indices at the census-tract level across the study cities.
\item Quantify the correlation between low MES and socioeconomic indicators:
median income, racial composition, vehicle ownership, and renter share.
\item Identify and rank the most mobility-disadvantaged tracts in each city,
distinguishing single-dimension from compounded multi-dimension deserts.
\item Build an interactive, open-access Plotly Dash dashboard for spatial
exploration and city comparison.
\item Produce a policy brief with actionable investment recommendations for city
transportation departments.
\end{enumerate}

\subsection{1.5\quad Significance}
Academically, this project fills a documented gap in the mobility-equity
literature by providing the first integrated four-dimension index at the
census-tract level with explicit, testable methodological choices. Practically,
the dashboard translates the analysis for non-technical stakeholders and supports
the federal Justice40 mandate to direct infrastructure benefits to disadvantaged
communities~\cite{argonne2022}, helping agencies target investment toward the
neighborhoods that need it most.

\subsection{1.6\quad Organization of the Report}
Chapter~2 reviews the literature on transportation equity, the four mobility
dimensions, composite-indicator construction, and spatial methods. Chapter~3
specifies the methodology in full. Chapters~4 and~5 (forthcoming) present results
and discussion.

% ============================== CHAPTER 2 ==============================
\section{Literature Review}

\subsection{2.1\quad Transit Equity and Transit Deserts}
Public transit accessibility is both a material and a social phenomenon shaping
employment access and quality of life~\cite{ermagun2020}. Aman and
Smith-Colin~\cite{aman2020} introduced a transit-desert framework quantifying the
supply--demand gap at the census-tract level using GTFS data, finding systematic
under-provision in low-income and minority tracts. Maharjan et al.~\cite{maharjan2024}
confirmed this nationally, showing transit-access gaps disproportionately affect
underserved populations across American metros. The federal Justice40 Initiative,
directing 40\% of infrastructure benefits to disadvantaged communities, demands
exactly the granular, neighborhood-level measurement tools this project
builds~\cite{argonne2022}.

\subsection{2.2\quad Urban Walkability and Equity}
The EPA National Walkability Index---incorporating intersection density,
land-use mix, and destination proximity---is the standard neighborhood
walkability measure~\cite{epa2024}. Research consistently shows walkability is
unevenly distributed, with communities of color inhabiting systematically less
walkable neighborhoods~\cite{urban2022}. Toolkits such as the 15-minute-city
score~\cite{albashir2024} operationalize walkable access but stop short of
integrating other mobility modes.

\subsection{2.3\quad Bicycle Infrastructure Equity}
Bicycle infrastructure is a recognized component of equitable, low-carbon
mobility, yet provision is uneven: dedicated facilities are concentrated in
higher-income areas, and bikeability indices remain
fragmented~\cite{ahmed2024,sustrans2023}. OpenStreetMap, accessed through the
OSMnx library~\cite{boeing2017}, has become the standard open source for
extracting and analyzing bicycle networks at scale.

\subsection{2.4\quad Electric-Vehicle Charging Equity}
As fleets electrify, equitable access to public charging---especially for
households without home charging---is an emerging equity concern. Lou et
al.~\cite{lou2025} analyzed 121 million U.S. households and found disadvantaged
communities have 64\% fewer public chargers per capita, a gap rising to 73\%
among renters. Chen and Li~\cite{chenli2024} found that 60--80\% of U.S. census
tracts have no public charging access at all, with extreme geographic
concentration.

\subsection{2.5\quad Composite Indicators in Public Policy}
Composite indicators compress multiple variables into a single, communicable
number; the OECD/JRC \emph{Handbook on Constructing Composite
Indicators}~\cite{nardo2008} is the canonical methodological reference,
structuring framework definition, normalization, weighting, and
robustness/sensitivity analysis. It stresses that normalization and weighting are
value-laden and must be paired with sensitivity analysis---a requirement this
study meets through its weighting-scheme and threshold tests (\S3.7--\S3.8).

\subsection{2.6\quad Spatial Analysis Methods}
Because mobility infrastructure is spatially structured, the analysis uses
established spatial-statistical tools. Global and local Moran's
$I$~\cite{anselin1995} test whether MES values cluster in space and locate
significant low-low cold spots. K-means clustering with silhouette-based model
selection~\cite{rousseeuw1987} derives an interpretable typology of neighborhood
mobility profiles.

\subsection{2.7\quad Gap in Existing Tools}
The EPA Smart Location Database, DOE AFDC locator, the Urban Institute SITE
tool~\cite{urbanSITE2025}, and the Oliver Wyman Forum Urban Mobility Readiness
Index~\cite{oliverwyman2024} each address one dimension or one geographic scale.
None combines all four mobility dimensions into a tract-level equity score.
Pandey et al.~\cite{pandey2025} found that rising urbanization is associated with
increasing infrastructure inequality, reinforcing the urgency of targeted,
evidence-based intervention tools.

% ============================== CHAPTER 3 ==============================
\section{Methodology}

\subsection{3.1\quad Study Area}
The unit of analysis is the U.S. Census tract (average 1,500--4,000 residents).
Three structurally distinct cities were selected for Phase~1 to stress-test the
method across mobility regimes (Table~\ref{tab:cities}); Phase~2 will extend to
Phoenix AZ, Atlanta GA, and Detroit MI after pipeline validation. Selection
criteria follow the proposal: population above 300,000, complete GTFS feeds,
demographic diversity, and variation in transit-investment history.

\begin{table}[h]
\centering
\caption{Phase-1 study cities and Census FIPS codes.}
\label{tab:cities}
\small
\begin{tabularx}{\textwidth}{l l l r Y}
\toprule
\textbf{City} & \textbf{State (FIPS)} & \textbf{County (FIPS)} & \textbf{Tracts} & \textbf{Rationale} \\
\midrule
Chicago & Illinois (17) & Cook (031) & 1{,}332 & Mature high-frequency transit; diverse \\
Houston & Texas (48) & Harris (201) & 1{,}115 & Auto-oriented, sprawling; large low-income pop. \\
Seattle & Washington (53) & King (033) & 495 & Progressive transit and EV investment \\
\bottomrule
\end{tabularx}
\end{table}

\subsection{3.2\quad Mobility Equity Score (MES)}
The MES is a composite index computed for each census tract from four
equally-weighted sub-indices (25\% each), each normalized to a 0--100 scale by
min--max scaling \emph{within each city} before aggregation. Higher scores
indicate greater mobility access relative to city peers. Table~\ref{tab:subindex}
defines each sub-index.

\begin{table}[h]
\centering
\caption{The four MES sub-indices.}
\label{tab:subindex}
\small
\begin{tabularx}{\textwidth}{l c Y l}
\toprule
\textbf{Sub-index} & \textbf{Weight} & \textbf{Key metric(s)} & \textbf{Data source} \\
\midrule
Transit access & 25\% & Stop density; AM-peak frequency (trips/hr); service span within 0.5 mi of tract centroid & GTFS (transit.land / agency) \\
Walkability & 25\% & EPA National Walkability Index (intersection density, land-use mix, destination proximity) & EPA Smart Location DB \\
Bicycle infrastructure & 25\% & Bike-lane km per km$^2$ of tract area & OpenStreetMap via OSMnx \\
EV charging access & 25\% & Public charger density per 1{,}000 residents; distance to nearest DC-fast station & DOE AFDC / OpenStreetMap \\
\bottomrule
\end{tabularx}
\end{table}

The composite is the weighted sum
$\mathrm{MES} = \sum_k w_k \cdot s_k^{\text{norm}}$ with $\sum_k w_k = 1$ and a
baseline of equal weights. Tracts scoring below the 25th percentile of the
composite MES are flagged as \emph{mobility deserts}; tracts below the 25th
percentile in two or more sub-indices simultaneously are \emph{compounded mobility
deserts}. The MES is cross-referenced with Census ACS variables: median household
income, percent non-white, percent zero-vehicle households, and percent
renter-occupied units.

\subsection{3.3\quad Data Sources}
All data are freely available with no licensing restrictions
(Table~\ref{tab:sources}).

\begin{table}[h]
\centering
\caption{Data sources and access.}
\label{tab:sources}
\small
\begin{tabularx}{\textwidth}{l Y l}
\toprule
\textbf{Source} & \textbf{Variables} & \textbf{Access} \\
\midrule
GTFS (transit.land / agency; Mobility Database) & Stop locations, schedules, headways & Free / token \\
EPA Smart Location Database v3 & National Walkability Index (block group) & Free download \\
OpenStreetMap / OSMnx & Bike lanes, paths; EV charging stations & Free Python library \\
DOE AFDC & EV charger locations, level, network & Free REST API (key) \\
US Census ACS 5-Year (2019--2023) & Income, race, vehicle ownership, tenure & Free Census API (key) \\
Census TIGER/Line (2023, 2019) & Tract \& block-group boundaries & Free download \\
\bottomrule
\end{tabularx}
\end{table}

\subsection{3.4\quad Data Collection and Preprocessing}
All layers are reprojected to EPSG:4326 for storage and mapping and to USA
Contiguous Albers Equal Area (EPSG:5070) for any length/area/distance
computation, so metrics are in true metres. GEOIDs are zero-padded strings (11
digits for tracts, 12 for block groups) to guarantee exact joins. A key
preprocessing problem is \emph{geographic-vintage mismatch}: the EPA Smart
Location Database is published on 2019-vintage block groups, whereas the study
uses 2023 tracts (to align with the 2019--2023 ACS). A naive GEOID-prefix join
loses tracts wherever boundaries were redrawn---about 57\% of Houston's. The
pipeline therefore performs \emph{areal interpolation}: the National Walkability
Index is joined to 2019 block-group polygons (a 100\% GEOID match) and
area-weighted onto the 2023 tracts, each tract's value being the area-weighted
mean of its overlapping block-group pieces. Missing values are preserved, not
imputed, so data gaps remain visible.

\subsection{3.5\quad Sub-Index Computation}
\textbf{Transit.} From GTFS \texttt{stops} and \texttt{stop\_times}, AM-peak
(07:00--09:00) departures per stop are converted to trips/hour and the daily
service span is computed. Stops are spatially joined to a 0.5-mile buffer of each
tract centroid, yielding stop density, mean frequency, and span, which are
normalized and averaged.
\textbf{Walkability.} Area-weighted National Walkability Index per tract (\S3.4),
normalized to 0--100.
\textbf{Bicycle.} OSM edges are retained as dedicated infrastructure when the
\texttt{highway} tag is cycleway/path/track or a positive \texttt{cycleway} tag
is present; clipped to tracts, bike-lane km per km$^2$ is normalized to 0--100.
\textbf{EV charging.} Public stations are spatially joined to tracts; chargers
per 1{,}000 residents and inverted distance to the nearest DC-fast station are
normalized and averaged. The authoritative source is the DOE/NREL AFDC; where
that API is unreachable, the pipeline falls back to OSM
\texttt{amenity=charging\_station} points, a key-free lower-bound proxy.

\subsection{3.6\quad Normalization and MES Construction}
Each raw sub-index is rescaled by min--max normalization,
$x_{\text{norm}} = \dfrac{x - x_{\min}}{x_{\max}-x_{\min}}\times 100$, computed
within each city, and combined as the equally-weighted sum in \S3.2. A robust
percentile-rank variant is retained for sensitivity checks.

\subsection{3.7\quad Weighting Methodology and Sensitivity Analysis}
Because equal weighting is itself a value choice, three alternatives are computed
and compared: \emph{entropy} weighting (objective weights from each sub-index's
dispersion), \emph{PCA} weighting (first-principal-component loadings), and
\emph{AHP} weighting (expert pairwise judgement). For each scheme the MES and the
equity correlations are recomputed, and the stability of the conclusions is
reported~\cite{nardo2008}.

\subsection{3.8\quad Mobility-Desert Classification}
A mobility desert is a tract below the city's 25th MES percentile; a compounded
desert is below the 25th percentile in $\geq 2$ sub-indices. Because the cut is
arbitrary, a threshold sensitivity analysis recomputes the desert set at the
20th, 25th, and 30th percentiles and reports the Jaccard overlap with the
baseline and the demographic profile of flagged tracts at each threshold.

\subsection{3.9\quad Equity Correlation Analysis}
Pearson correlations relate the MES (and each sub-index) to the four demographic
variables, with two-sided $p$-values and 95\% confidence intervals from the
Fisher $z$-transformation, $z=\operatorname{arctanh}(r)$,
$\mathrm{SE}=1/\sqrt{n-3}$, $\mathrm{CI}=\tanh(z\pm 1.96\,\mathrm{SE})$. A negative
MES--income or MES--percent-non-white correlation would indicate that
more-disadvantaged neighborhoods receive less mobility infrastructure.

\subsection{3.10\quad Spatial Autocorrelation}
Global Moran's $I$ tests spatial clustering of MES using queen-contiguity,
row-standardized weights with 999 permutations; Local Moran's $I$ (LISA)
classifies each tract as high-high, low-low, high-low, low-high, or
non-significant, mapping low-low clusters as mobility cold spots~\cite{anselin1995}.

\subsection{3.11\quad Cluster Analysis}
K-means clustering on the four standardized sub-indices derives a typology of
neighborhood mobility profiles; $k$ is chosen by maximizing the mean silhouette
score over $k\in\{2,\dots,7\}$, with the inertia (elbow) curve reported for
validation~\cite{rousseeuw1987}. Each cluster is labelled from its centroid
profile.

\subsection{3.12\quad Dashboard Development}
The dashboard (Plotly Dash with Dash Bootstrap Components) provides a tract-level
choropleth of the MES or any sub-index with desert filters; a click-through side
panel (score, city percentile, sub-index bars versus the city median,
auto-generated interpretation, and demographics); a city equity-gap comparison
(median MES of the poorest versus richest income quintile); a ranked table of the
most disadvantaged tracts; and CSV/PNG export.

\subsection{3.13\quad Reproducibility and Implementation Status}
All processing is in Python~3.11 (pandas, geopandas, numpy, scipy, scikit-learn,
PySAL). Every stage is an independently-runnable module, raw downloads are cached,
and a \texttt{pytest} suite covers normalization, MES construction, and the
correlation statistics. At submission the pipeline runs end-to-end for all three
Phase-1 cities on measured open data across every dimension (transit via GTFS for
all three cities, including METRO Houston resolved through the Mobility Database
API). Global Moran's $I$ for the MES is 0.82--0.92 ($p\approx 0.001$) across the
three cities, confirming strong spatial clustering of mobility access; the
substantive equity findings are deferred to Chapter~4.

% ============================== REFERENCES ==============================
\section*{References}
\addcontentsline{toc}{section}{References}
\small
\begin{thebibliography}{99}
\setlength{\itemsep}{1pt}

\bibitem{aman2020} J. J. Aman and J. Smith-Colin, ``Transit deserts: Equity analysis of public transit accessibility,'' \textit{Journal of Transport Geography}, vol.~89, p.~102869, 2020.

\bibitem{urban2022} Urban Institute, ``Redefining Walkability: A Neighborhood-Level Analysis,'' Urban Institute, Washington, DC, 2022. [Online]. Available: \url{https://www.urban.org/data-tools/redefining-walkability}

\bibitem{lou2025} J. Lou, A. Chandra, and M. Gordon, ``Equity and reliability of public electric vehicle charging stations in the United States,'' \textit{Nature Communications}, vol.~16, p.~4201, 2025.

\bibitem{ahmed2024} S. Ahmed, Z. Sultana, and M. Rahman, ``A systematic review of bicycle infrastructure design principles within urban bikeability indices,'' \textit{Sustainability}, vol.~17, no.~14, p.~6409, 2024.

\bibitem{ermagun2020} A. Ermagun and N. Tilahun, ``Transit equity and job accessibility,'' \textit{Transportation Research Part A}, vol.~134, pp.~191--204, 2020.

\bibitem{maharjan2024} S. Maharjan, F. Janatabadi, and A. Ermagun, ``Spatial inequity of transit and automobile access gap across America for underserved populations,'' \textit{Transportation Research Record}, vol.~2678, no.~1, pp.~674--690, 2024.

\bibitem{argonne2022} Argonne National Laboratory, ``Electric Vehicle Charging Equity Considerations,'' U.S. Department of Energy, 2022. [Online]. Available: \url{https://www.anl.gov/esia}

\bibitem{epa2024} U.S. EPA, ``Smart Location Mapping and National Walkability Index,'' EPA, 2024. [Online]. Available: \url{https://www.epa.gov/smartgrowth/smart-location-mapping}

\bibitem{sustrans2023} Sustrans, ``Cycling and Walking Investment Strategy: Equity in Active Travel Infrastructure,'' Sustrans, Bristol, UK, 2023.

\bibitem{albashir2024} A. Albashir et al., ``15min City Score Toolkit: Urban Walkability Analytics,'' Transform Transport Working Papers, Zenodo, 2024.

\bibitem{chenli2024} X. Chen and Y. Li, ``Electric vehicle charging equity and accessibility: A U.S. policy analysis,'' \textit{Transportation Research Part D}, vol.~118, p.~103748, 2024.

\bibitem{urbanSITE2025} Urban Institute, ``Spending on Infrastructure toward Equity (SITE) Tool,'' 2025. [Online]. Available: \url{https://www.urban.org}

\bibitem{oliverwyman2024} Oliver Wyman Forum, ``2024 Urban Mobility Readiness Index,'' 2024. [Online]. Available: \url{https://www.oliverwymanforum.com/mobility}

\bibitem{pandey2025} B. Pandey, C. Brelsford, and K. C. Seto, ``Rising infrastructure inequalities accompany urbanization and economic development,'' \textit{Nature Communications}, vol.~16, p.~1082, 2025.

\bibitem{boeing2017} G. Boeing, ``OSMnx: New methods for acquiring, constructing, analyzing, and visualizing complex street networks,'' \textit{Computers, Environment and Urban Systems}, vol.~65, pp.~126--139, 2017.

\bibitem{nardo2008} M. Nardo, M. Saisana, A. Saltelli, et al., \textit{Handbook on Constructing Composite Indicators: Methodology and User Guide}. OECD/JRC, 2008.

\bibitem{anselin1995} L. Anselin, ``Local Indicators of Spatial Association---LISA,'' \textit{Geographical Analysis}, vol.~27, no.~2, pp.~93--115, 1995.

\bibitem{rousseeuw1987} P. J. Rousseeuw, ``Silhouettes: A graphical aid to the interpretation and validation of cluster analysis,'' \textit{Journal of Computational and Applied Mathematics}, vol.~20, pp.~53--65, 1987.

\end{thebibliography}

\end{document}
